\documentclass[10pt,a4paper]{amsart}
\usepackage{amsmath,amssymb,amsthm,graphicx}
\usepackage{verbatim}
\newenvironment{reference}{\comment}{\endcomment}
%\newenvironment{reference}{}{}
\newenvironment{slogan}{\comment}{\endcomment}
\newenvironment{history}{\comment}{\endcomment}

% For commutative diagrams we use Xy-pic
\usepackage[all]{xy}

% We use 2cell for 2-commutative diagrams.
\xyoption{2cell}
% Theorem environments.
%
\theoremstyle{plain}
\newtheorem{theorem}[subsection]{Theorem}
\newtheorem{proposition}[subsection]{Proposition}
\newtheorem{lemma}[subsection]{Lemma}

\theoremstyle{definition}
\newtheorem{definition}[subsection]{Definition}
\newtheorem{example}[subsection]{Example}
\newtheorem{exercise}[subsection]{Exercise}
\newtheorem{situation}[subsection]{Situation}

\theoremstyle{remark}
\newtheorem{remark}[subsection]{Remark}
\newtheorem{remarks}[subsection]{Remarks}

\numberwithin{equation}{subsection}

% Macros
%
\def\lim{\mathop{\mathrm{lim}}\nolimits}
\def\colim{\mathop{\mathrm{colim}}\nolimits}
\def\Spec{\mathop{\mathrm{Spec}}}
\def\Hom{\mathop{\mathrm{Hom}}\nolimits}
\def\Ext{\mathop{\mathrm{Ext}}\nolimits}
\def\SheafHom{\mathop{\mathcal{H}\!\mathit{om}}\nolimits}
\def\SheafExt{\mathop{\mathcal{E}\!\mathit{xt}}\nolimits}
\def\Sch{\mathit{Sch}}
\def\Mor{\mathop{\mathrm{Mor}}\nolimits}
\def\Ob{\mathop{\mathrm{Ob}}\nolimits}
\def\Sh{\mathop{\mathit{Sh}}\nolimits}
\def\NL{\mathop{N\!L}\nolimits}
\def\CH{\mathop{\mathrm{CH}}\nolimits}
\def\proetale{{pro\text{-}\acute{e}tale}}
\def\etale{{\acute{e}tale}}
\def\QCoh{\mathit{QCoh}}
\def\Ker{\mathop{\mathrm{Ker}}}
\def\Im{\mathop{\mathrm{Im}}}
\def\Coker{\mathop{\mathrm{Coker}}}
\def\Coim{\mathop{\mathrm{Coim}}}

% Boxtimes
%
\DeclareMathSymbol{\boxtimes}{\mathbin}{AMSa}{"02}

%
% Macros for moduli stacks/spaces
%
\def\QCohstack{\mathcal{QC}\!\mathit{oh}}
\def\Cohstack{\mathcal{C}\!\mathit{oh}}
\def\Spacesstack{\mathcal{S}\!\mathit{paces}}
\def\Quotfunctor{\mathrm{Quot}}
\def\Hilbfunctor{\mathrm{Hilb}}
\def\Curvesstack{\mathcal{C}\!\mathit{urves}}
\def\Polarizedstack{\mathcal{P}\!\mathit{olarized}}
\def\Complexesstack{\mathcal{C}\!\mathit{omplexes}}
% \Pic is the operator that assigns to X its picard group, usage \Pic(X)
% \Picardstack_{X/B} denotes the Picard stack of X over B
% \Picardfunctor_{X/B} denotes the Picard functor of X over B
\def\Pic{\mathop{\mathrm{Pic}}\nolimits}
\def\Picardstack{\mathcal{P}\!\mathit{ic}}
\def\Picardfunctor{\mathrm{Pic}}
\def\Deformationcategory{\mathcal{D}\!\mathit{ef}}
\usepackage{cleveref}
\crefname{originalreference}{original reference}{original references}
\Crefname{originalreference}{Original reference}{Original references}
\makeatletter
\newlabel{lemma-count-embeddings-explicitly}{{09H8}{source}{}{sourcepointer-1}{}}
\newlabel{lemma-count-embeddings-explicitly@cref}{{[originalreference][0][]09H8}{[1][0][]source}{}{}{}}
\makeatother
\crefname{originalreference}{original reference}{original references}
\Crefname{originalreference}{Original reference}{Original references}
\makeatletter
\newlabel{lemma-finite-purely-inseparable}{{09HI}{source}{}{sourcepointer-2}{}}
\newlabel{lemma-finite-purely-inseparable@cref}{{[originalreference][0][]09HI}{[1][0][]source}{}{}{}}
\makeatother
\crefname{originalreference}{original reference}{original references}
\Crefname{originalreference}{Original reference}{Original references}
\makeatletter
\newlabel{lemma-separable-equality}{{09HA}{source}{}{sourcepointer-3}{}}
\newlabel{lemma-separable-equality@cref}{{[originalreference][0][]09HA}{[1][0][]source}{}{}{}}
\makeatother
\crefname{originalreference}{original reference}{original references}
\Crefname{originalreference}{Original reference}{Original references}
\makeatletter
\newlabel{lemma-separable-first}{{030K}{source}{}{sourcepointer-4}{}}
\newlabel{lemma-separable-first@cref}{{[originalreference][0][]030K}{[1][0][]source}{}{}{}}
\makeatother
\title{The separable degree equals the number of embeddings into an algebraic closure}
\author{The Stacks Project Authors}
\date{stacks-project:fields.tex:local-snapshot-20261012; GNU FDL 1.2 or later}
\begin{document}
\maketitle
\noindent\textit{Editorial note.} All delivered mathematical text is the author's own, copied from the pinned Stacks Project chapter file fields.tex; no statement or proof step was rewritten, reordered or removed.\par\medskip
\noindent\textit{Editorial note.} The counted proof is the author's complete argument: The unique extension claim ('by the result of the previous paragraph') is the just-proved purely inseparable case applied to K/K\_s: there is exactly one K-embedding of K into F-bar over a given embedding of K\_s.\par\medskip
\noindent\textit{Editorial note.} The label lemma-separable-degree names the counted statement, and the proof environment immediately following it in the source is delivered verbatim as its proof.\par\medskip
\noindent\textit{Editorial note.} The counted statement is carried by the author's own lemma environment in the chapter source.\par\medskip
\noindent\textit{Allowed standard prerequisites (editorial).}\begin{itemize}
\item Context span 1943-1954: the definition of the separable degree as the degree of the separable subextension
\item Same chapter: lemma-finite-purely-inseparable (the x\textasciicircum{}p - alpha\_i\textasciicircum{}p tower for purely inseparable extensions), lemma-count-embeddings-explicitly (|Mor\_F(K, F-bar)| = product deg\_s(P\_i)), lemma-separable-first (the separable subextension K\_s), lemma-separable-equality (|Mor\_F(K\_s, F-bar)| = [K\_s : F] with equality iff separable)
\end{itemize}\medskip
\ifdefined\hypertarget\hypertarget{sourcepointer-1}{}\fi
\noindent\textit{Source reference.} Original reference 09H8: sources/stacks-fields/pinned-original-source.tex, source lines 1543--1543. Prerequisite pointer invoked at 1979; it names a standard fact of the same chapter that the counted argument uses, and every step of the counted proof built on it is retained verbatim in the delivered body.\par\medskip
\ifdefined\hypertarget\hypertarget{sourcepointer-2}{}\fi
\noindent\textit{Source reference.} Original reference 09HI: sources/stacks-fields/pinned-original-source.tex, source lines 1886--1886. Prerequisite pointer invoked at 1977; it names a standard fact of the same chapter that the counted argument uses, and every step of the counted proof built on it is retained verbatim in the delivered body.\par\medskip
\ifdefined\hypertarget\hypertarget{sourcepointer-3}{}\fi
\noindent\textit{Source reference.} Original reference 09HA: sources/stacks-fields/pinned-original-source.tex, source lines 1601--1601. Prerequisite pointer invoked at 1986; it names a standard fact of the same chapter that the counted argument uses, and every step of the counted proof built on it is retained verbatim in the delivered body.\par\medskip
\ifdefined\hypertarget\hypertarget{sourcepointer-4}{}\fi
\noindent\textit{Source reference.} Original reference 030K: sources/stacks-fields/pinned-original-source.tex, source lines 1921--1921. Prerequisite pointer invoked at 1985; it names a standard fact of the same chapter that the counted argument uses, and every step of the counted proof built on it is retained verbatim in the delivered body.\par\medskip
\begin{definition}
\label{definition-insep-degree}
Let $E/F$ be an algebraic field extension. Let $E_{sep}$ be the subextension
found in Lemma \ref{lemma-separable-first}.
\begin{enumerate}
\item The integer $[E_{sep} : F]$ is called the {\it separable
degree} of the extension. Notation $[E : F]_s$.
\item The integer $[E : E_{sep}]$ is called the {\it inseparable
degree}, or the {\it degree of inseparability} of the extension.
Notation $[E : F]_i$.
\end{enumerate}
\end{definition}

\begin{lemma}
\label{lemma-separable-degree}
Let $K/F$ be a finite extension. Let $\overline{F}$ be an algebraic
closure of $F$. Then $[K : F]_s = |\Mor_F(K, \overline{F})|$.
\end{lemma}

\begin{proof}
We first prove this when $K/F$ is purely inseparable. Namely, we claim that
in this case there is a unique map $K \to \overline{F}$. This can be
seen by choosing a sequence of elements $\alpha_1, \ldots, \alpha_n \in K$
as in Lemma \ref{lemma-finite-purely-inseparable}. The irreducible polynomial
of $\alpha_i$ over $F(\alpha_1, \ldots, \alpha_{i - 1})$ is $x^p - \alpha_i^p$.
Applying Lemma \ref{lemma-count-embeddings-explicitly} we see that
$|\Mor_F(K, \overline{F})| = 1$. On the other hand, $[K : F]_s = 1$
in this case hence the equality holds.

\medskip\noindent
Let's return to a general finite extension $K/F$. In this case
choose $F \subset K_s \subset K$ as in Lemma \ref{lemma-separable-first}.
By Lemma \ref{lemma-separable-equality} we have
$|\Mor_F(K_s, \overline{F})| = [K_s : F] = [K : F]_s$.
On the other hand, every field map $\sigma' : K_s \to \overline{F}$
extends to a unique field map $\sigma : K \to \overline{F}$ by the
result of the previous paragraph. In other words
$|\Mor_F(K, \overline{F})| = |\Mor_F(K_s, \overline{F})|$
and the proof is done.
\end{proof}

\end{document}
